\documentclass[11pt,a4paper]{article}

\usepackage{geometry}
\geometry{a4paper, top=25mm, bottom=25mm, left=20mm, right=20mm}

\usepackage{xcolor}
\usepackage{booktabs}
\usepackage{tabularx}
\usepackage{fancyhdr}
\usepackage{enumitem}
\usepackage{hyperref}

\hypersetup{
    colorlinks=true,
    linkcolor=blue,
    citecolor=blue,
    urlcolor=blue,
    pdftitle={راهنمای جامع مهاجرت و انتقال سرور پلتفرم آموزشی لایتنر به همراه رمزها و ساختار SQL},
    pdfauthor={تیم توسعه سکوی یادگیری لایتنر}
}

% Header and Footer styling
\pagestyle{fancy}
\fancyhf{}
\fancyhead[R]{\small راهنمای جامع مهاجرت سرور و مستندات دیتابیس}
\fancyhead[L]{\small سکوی یادگیری لایتنر}
\fancyfoot[C]{\thepage}
\renewcommand{\headrulewidth}{0.4pt}
\renewcommand{\footrulewidth}{0.4pt}

% XePersian MUST be loaded as the last package
\usepackage{xepersian}
\settextfont{Vazirmatn}

\begin{document}

% ----------------------------------------------------------------------
% Cover / Title Section
% ----------------------------------------------------------------------
\begin{titlepage}
    \centering
    \vspace*{1.5cm}
    
    {\Huge\textbf{سکوی یادگیری لایتنر (RightLearn)}}\\[0.8cm]
    {\LARGE\textbf{سند فنی و دستورالعمل جامع انتقال و مهاجرت سرور}}\\[0.4cm]
    {\Large به همراه کلیه رمزهای عبور دیتابیس، مشخصات امنیتی و فایل‌های مایگریشن \lr{SQL}}\\[1.5cm]
    
    \vspace*{2cm}
    
    \begin{trivlist}
        \item \textbf{خلاصه اهداف و محتوای سند:}
        \item این مستند حاوی کلیه اطلاعات حساس و فنی مورد نیاز برای انتقال، پشتیبان‌گیری، بازیابی و مدیریت زیرساخت سرور سکوی یادگیری لایتنر می‌باشد. تمامی رمزهای عبور دیتابیس \verb|PostgreSQL|، ردیس، کلیدهای احراز هویت \verb|JWT|، درگاه پیامک و پرداخت به همراه شرح دقیق ۱۷ فایل مایگریشن \verb|SQL| در این سند گنجانده شده است تا سرور جدید با تضمین عدم از دست رفتن اطلاعات (\lr{Zero Data Loss}) راه‌اندازی گردد.
    \end{trivlist}
    
    \vfill
    {\small این سند حاوی اطلاعات محرمانه سیستم است و باید در نگهداری آن دقت لازم به عمل آید.}
\end{titlepage}

\tableofcontents
\newpage

% ----------------------------------------------------------------------
% Section 1
% ----------------------------------------------------------------------
\section{مقدمه و اهداف مهاجرت}
هنگامی که نیاز به ارتقای سخت‌افزار، تغییر دیتاسنتر یا جابه‌جایی سرور میزبان سکوی لایتنر ایجاد می‌شود، عملیات انتقال باید با رعایت کامل پروتکل‌های امنیتی و با هدف \textbf{عدم از دست رفتن حتی یک رکورد اطلاعاتی} (\lr{Zero Data Loss}) و \textbf{کوتاه‌ترین زمان قطعی ممکن} (\lr{Minimal Downtime}) صورت پذیرد.

این راهنما شامل کلیه مراحل از آماده‌سازی دامنه، کانفیگ لینوکس اوبونتو، نصب ابزارهای کانتینری داکر، پشتیبان‌گیری و بازیابی پایگاه داده \verb|PostgreSQL|، انتقال داده‌های پایدار، راه‌اندازی پروکسی معکوس \verb|Nginx|، صدور مجدد گواهی امنیتی \verb|SSL|، تحویل فایل‌های اسکریپت دیتابیس \verb|SQL| و بازنشانی اعتبارسنجی‌ها می‌باشد.

% ----------------------------------------------------------------------
% Section 2
% ----------------------------------------------------------------------
\section{معماری سرویس‌ها و داده‌های مستقر در سرور}
ساختار اجرایی سامانه مبتنی بر معماری میکروسرویس کانتینری داکر (\verb|Docker Compose|) به همراه لایه پروکسی وب در سطح سیستم‌عامل میزبان می‌باشد:

\subsection{سرویس‌های کانتینری (\lr{Docker Stack})}
\begin{itemize}
    \item \textbf{پایگاه داده (\lr{leitner-postgres-db}):} دیتابیس \verb|PostgreSQL 16 Alpine| حاوی اطلاعات کاربران، سوابق لایتنر، لاگ‌ها، سفارشات، پرداخت‌ها و تنظیمات سیستم.
    \item \textbf{سرویس کش و صف پیام (\lr{leitner-redis-cache}):} سرور \verb|Redis 7.2 Alpine| جهت نشست‌ها، کش داده‌ها و رویدادهای بی‌درنگ.
    \item \textbf{هسته سرویس‌دهنده و وب‌سرویس (\lr{leitner-backend-api}):} سرویس \verb|ASP.NET Core 8.0| پاسخگوی درخواست‌های اپلیکیشن موبایل و پنل وب روی پورت داخلی \verb|8080|.
    \item \textbf{پنل وب مدیریت (\lr{leitner-admin-panel}):} واسط کاربری تحت وب مبتنی بر \verb|React/Vite| سرویس‌دهنده روی پورت داخلی \verb|8081| (یا \verb|80|).
    \item \textbf{پردازشگر وظایف پس‌زمینه (\lr{leitner-background-worker}):} سرویس زمان‌بندی و پشتیبان‌گیری خودکار مبتنی بر دات‌نت و \verb|Hangfire|.
\end{itemize}

\subsection{داده‌ها و فایل‌های پایدار که باید انتقال یابند}
\begin{enumerate}
    \item \textbf{دامپ دیتابیس \lr{leitner\_db}:} تمامی جداول، داده‌ها و ایندکس‌های پایگاه داده.
    \item \textbf{فایل‌های فشرده دوره‌ها (\lr{Course Packages}):} بسته‌های دوره‌ها شامل فایل‌های دیتابیس آفلاین \verb|.db|، پوشه‌های صوتی و تصاویر دوره‌ها در مسیر:\\
    \verb|/opt/leitner-platform/backend/LeitnerPlatform.API/wwwroot/courses/|
    \item \textbf{تصاویر آپلودی کاربران و برندسازی:} آواتارها و لوگوی سامانه در مسیر:\\
    \verb|/opt/leitner-platform/backend/LeitnerPlatform.API/wwwroot/uploads/|
    \item \textbf{فایل متغیرهای محیطی (\lr{.env}):} شامل کلیه رمزهای عبور، کلیدهای دسترسی پیامک و درگاه.
\end{enumerate}

\begin{table}[ht]
\centering
\caption{جدول مشخصات پورت‌ها و پروتکل‌های سرویس‌های سرور}
\vspace{2mm}
\begin{tabularx}{\textwidth}{llllX}
\toprule
\textbf{نام سرویس} & \textbf{پورت هاست} & \textbf{پورت کانتینر} & \textbf{پروتکل} & \textbf{توضیحات و کاربرد} \\
\midrule
Nginx (Host) & 80, 443 & - & HTTP/HTTPS & پروکسی وب و ورودی اصلی \\
Web API & 8080 & 8080 & HTTP & هسته وب‌سرویس بک‌اند \\
Admin Web & 8081 & 80 & HTTP & پنل مدیریت لایتنر \\
PostgreSQL & 5432 & 5432 & TCP & پایگاه داده لایتنر \\
Redis & 6379 & 6379 & TCP & کش و صف رویداد \\
\bottomrule
\end{tabularx}
\end{table}

% ----------------------------------------------------------------------
% Section 3: Passwords and Credentials
% ----------------------------------------------------------------------
\newpage
\section{اطلاعات محرمانه، رمزهای عبور دیتابیس و اعتبارسنجی‌ها}
در این بخش کلیه کلمات عبور، نام‌های کاربری و رشته‌های اتصال به پایگاه داده و سرویس‌ها به طور کامل مستند شده است:

\subsection{اطلاعات دسترسی به پایگاه داده \lr{PostgreSQL 16}}
\begin{table}[ht]
\centering
\caption{مشخصات اتصال و کلمات عبور پایگاه داده اصلی}
\vspace{2mm}
\begin{tabularx}{\textwidth}{lX}
\toprule
\textbf{پارامتر} & \textbf{مقدار و مقداردهی محیطی} \\
\midrule
نام پایگاه داده (\lr{Database Name}) & \verb|leitner_db| \\
نام کاربری مدیر دیتابیس (\lr{Database User}) & \verb|leitner_admin| \\
رمز عبور دیتابیس (\lr{Database Password}) & \verb|leitner_secure_pass_2026| \\
پورت پیش‌فرض کانتینر و هاست & \verb|5432| \\
میزبان درون شبکه داکر (\lr{Internal Host}) & \verb|db| (یا \verb|leitner-postgres-db|) \\
میزبان خارج از کانتینر (\lr{Host/VPS}) & \verb|localhost| یا \verb|127.0.0.1| \\
\bottomrule
\end{tabularx}
\end{table}

\textbf{رشته اتصال درون شبکه کانتینری داکر:}
\begin{latin}
\footnotesize
\begin{verbatim}
Server=db;Port=5432;Database=leitner_db;User Id=leitner_admin;Password=leitner_secure_pass_2026;
\end{verbatim}
\end{latin}

\textbf{رشته اتصال از طریق هاست یا نرم‌افزارهای مدیریت دیتابیس (مانند DBeaver یا pgAdmin):}
\begin{latin}
\footnotesize
\begin{verbatim}
Server=127.0.0.1;Port=5432;Database=leitner_db;User Id=leitner_admin;Password=leitner_secure_pass_2026;
\end{verbatim}
\end{latin}

\subsection{اطلاعات اتصال سرویس کش و صف پیام \lr{Redis 7.2}}
\begin{table}[ht]
\centering
\caption{مشخصات اتصال و رمز عبور سرویس ردیس}
\vspace{2mm}
\begin{tabularx}{\textwidth}{lX}
\toprule
\textbf{پارامتر} & \textbf{مقدار} \\
\midrule
میزبان کانتینری (\lr{Host}) & \lr{\texttt{redis}} (یا \lr{\texttt{leitner-redis-cache}}) \\
پورت اتصال (\lr{Port}) & \verb|6379| \\
رمز عبور عبور امنیتی (\lr{Requirepass}) & \lr{\texttt{redis\_secure\_pass\_2026}} \\
رشته اتصال در هسته دات‌نت (\lr{Connection String}) & {\small\lr{\texttt{redis:6379,password=redis\_secure\_pass\_2026}}} \\
\bottomrule
\end{tabularx}
\end{table}

\subsection{کلیدهای امنیتی احراز هویت \lr{JWT} و توکن‌ها}
\begin{itemize}
    \item \textbf{کلید امضای رمزنگاری توکن (\lr{JWT\_SECRET\_KEY}):}
    \begin{latin}
    \lr{\texttt{jwt\_secret\_lts\_2026\_super\_secure\_key}}
    \end{latin}
    \textbf{نکته مهم:} این کلید در سرور جدید دقیقاً باید با سرور قبلی یکسان باشد؛ در غیر این صورت، نشست ورود تمام کاربرانی که در اپلیکیشن لاگین کرده‌اند منقضی شده و ناچار به لاگین مجدد می‌شوند.
    \item \textbf{صادرکننده توکن (\lr{Issuer}) و مخاطب (\lr{Audience}):} \verb|LeitnerLearningPlatform|
    \item \textbf{مدت اعتبار توکن دسترسی (\lr{Access Token}):} ۱۵ دقیقه (تجدید خودکار با Refresh Token تا ۳۰ روز).
\end{itemize}

\subsection{اطلاعات درگاه پیامک و پرداخت ریالی}
\begin{table}[ht]
\centering
\caption{کلیدهای درگاه پیامک و درگاه پرداخت آنلاین}
\vspace{2mm}
\begin{tabularx}{\textwidth}{lp{10.5cm}}
\toprule
\textbf{سرویس} & \textbf{کلید و تنظیمات ثبت‌شده} \\
\midrule
ارائه‌دهنده پیامک (\lr{SMS Provider}) & \lr{\texttt{IranPayamak}} (و سازگار با \lr{\texttt{Kavenegar}}) \\
کلید وب‌سرویس پیامک (\lr{API Key}) & {\footnotesize\lr{\texttt{AfhU2pzBoaWmXI3HYx7YipKbQBGaBYBYjpRSMLTSO1bOSICWeV}}} \\
سرشماره اختصاصی ارسال پیامک & \verb|50002178584000| \\
کد الگوی اعتبارسنجی (\lr{Pattern Code}) & \verb|coHfKHvJK6| \\
مرچنت درگاه پرداخت زرین‌پال (\lr{Merchant ID}) & \lr{\texttt{167ccd1b-f5d3-407d-85d1-a73c4f2ba3eb}} \\
آدرس کال‌بک بازگشت تراکنش & {\footnotesize\url{https://api.rightlearn.ir/api/v1/purchases/zarinpal/callback}} \\
\bottomrule
\end{tabularx}
\end{table}

\subsection{اطلاعات دسترسی به سرور لینوکس میزبان (\lr{VPS})}
\begin{itemize}
    \item \textbf{نام کاربری سیستمی:} \verb|root|
    \item \textbf{آی‌پی سرور فعلی مبدا:} \verb|45.94.215.188|
    \item \textbf{رمز عبور کاربر root سرور فعلی:} \verb|UuP>wKyk45h,bw4b2026|
    \item \textbf{کلید احراز هویت بدون رمز (SSH Key):} کلید خصوصی \verb|id_rsa_deploy|
\end{itemize}

\subsection{محتوای استاندارد فایل پیکربندی \lr{.env} سرور}
فایل \verb|/opt/leitner-platform/.env| در سرور جدید باید دارای مقادیر زیر باشد:

\begin{latin}
\small
\begin{verbatim}
# Database & Cache Configurations
DB_PASSWORD=leitner_secure_pass_2026
REDIS_PASSWORD=redis_secure_pass_2026

# Backend Security & JWT
ASPNETCORE_ENVIRONMENT=Production
JWT_SECRET_KEY=jwt_secret_lts_2026_super_secure_key

# Ports
ADMIN_PORT=8081
BACKEND_PORT=8080

# SMS OTP Provider
SMS_PROVIDER=IranPayamak
SMS_GATEWAY_API_KEY=AfhU2pzBoaWmXI3HYx7YipKbQBGaBYBYjpRSMLTSO1bOSICWeV
SMS_SENDER=50002178584000
SMS_PATTERN_CODE=coHfKHvJK6

# ZarinPal IPG
ZARINPAL_MERCHANT_ID=167ccd1b-f5d3-407d-85d1-a73c4f2ba3eb
ZARINPAL_SANDBOX=false
ZARINPAL_CALLBACK_URL=https://api.rightlearn.ir/api/v1/purchases/zarinpal/callback
\end{verbatim}
\end{latin}

% ----------------------------------------------------------------------
% Section 4: SQL Migration Files
% ----------------------------------------------------------------------
\newpage
\section{ساختار پایگاه داده و فایل‌های اسکریپت مایگریشن \lr{SQL}}
سکوی لایتنر از موتور مایگریشن استاندارد \lr{DbUp} استفاده می‌کند. تمامی اسکریپت‌های ساختار پایگاه داده در پوشه پروژه به نشانی زیر قرار دارند:
\begin{latin}
\verb|deployment/db/migrations/|
\end{latin}
هنگامی که کانتینر بک‌اند لایتنر روشن می‌شود، به صورت خودکار فایل‌های مایگریشن را به ترتیب نسخه بررسی کرده و در صورت عدم اجرا، آن‌ها را روی پایگاه داده اعمال می‌نماید. وضعیت اسکریپت‌های اجرا شده در جدول \verb|schemaversions| دیتابیس ذخیره می‌شود.

\subsection{فهرست کامل ۱۷ اسکریپت مایگریشن \lr{SQL} سرور}
تمامی این فایل‌ها در مسیر ذکر شده موجود بوده و وظیفه ایجاد جداول و تکامل ساختار اطلاعاتی سیستم را بر عهده دارند:

\begin{itemize}[leftmargin=0.8cm]
    \item \textbf{\lr{\texttt{V1\_\_Initial\_Schema.sql}}}:\\
    ساخت جداول پایه هسته سیستم شامل \lr{\texttt{users}}، \lr{\texttt{courses}}، \lr{\texttt{cards}}، \lr{\texttt{leitner\_progress}}، \lr{\texttt{purchases}}، \lr{\texttt{flashcard\_reports}} و \lr{\texttt{banners}}. تولید خودکار شناسه‌های یکتا (\lr{UUID}) و ایجاد ایندکس‌های کلیدی شماره موبایل و دسته‌بندی دوره‌ها.

    \item \textbf{\lr{\texttt{V2\_\_Add\_User\_IsAdmin.sql}}}:\\
    افزودن ستون \lr{\texttt{is\_admin}} از نوع بولین به جدول کاربران برای تفکیک دسترسی مدیران سیستم از کاربران عادی.

    \item \textbf{\lr{\texttt{V3\_\_Add\_Course\_Fields.sql}}}:\\
    افزودن ستون‌های مدت تخمینی دوره (\lr{\texttt{estimated\_duration}}) و تعداد کل کارت‌ها به جدول دوره‌ها.

    \item \textbf{\lr{\texttt{V4\_\_Add\_Remote\_Config.sql}}}:\\
    ایجاد جدول کلید-مقدار تنظیمات زنده سیستم (\lr{\texttt{system\_configs}}) جهت مدیریت وضعیت حالت تعمیرات (\lr{\texttt{maintenance\_mode}}) و هدر اطلاعیه‌ها.

    \item \textbf{\lr{\texttt{V5\_\_Seed\_Sample\_Courses.sql}}}:\\
    درج داده‌های اولیه و نمونه دوره‌های آموزشی زبان انگلیسی و پرسش‌های فلش‌کارت جهت تست و اعتبارسنجی اولیه.

    \item \textbf{\lr{\texttt{V6\_\_Update\_Course\_Download\_Urls.sql}}}:\\
    تنظیم و یکپارچه‌سازی لینک‌های دانلود مستقیم بسته‌های فشرده دوره‌ها (\lr{\texttt{.zip}}) از مسیر وب‌سرور.

    \item \textbf{\lr{\texttt{V7\_\_Add\_Report\_UserMobileNumber.sql}}}:\\
    ثبت خودکار شماره همراه گزارش‌دهنده خطا در جدول گزارش‌های فلش‌کارت جهت تماس و پشتیبانی آموزشی.

    \item \textbf{\lr{\texttt{V8\_\_Course\_Soft\_Delete\_And\_Versioning.sql}}}:\\
    پیاده‌سازی مکانیزم حذف نرم (\lr{\texttt{is\_deleted}}) و نگهداری شماره نسخه به‌روزرسانی محتوای دوره‌ها.

    \item \textbf{\lr{\texttt{V9\_\_Seed\_1100\_Words\_Course.sql}}}:\\
    ثبت مشخصات، تصاویر و متادیتای رسمی دوره پرطرفدار ۱۱۰۰ واژه ضروری بارونز در کاتالوگ دوره‌ها.

    \item \textbf{\lr{\texttt{V10\_\_Add\_Social\_And\_Support\_Configs.sql}}}:\\
    درج مشخصات کانال روبیکا، تلگرام، شماره تلفن پشتیبانی و ایمیل ارتباطی در جدول تنظیمات ریموت.

    \item \textbf{\lr{\texttt{V11\_\_Add\_Leitner\_Interval\_Configs.sql}}}:\\
    تعریف فواصل زمانی روزهای جعبه لایتنر (روزهای ۱، ۲، ۴، ۸، ۱۶، ۳۰) در جدول \lr{\texttt{system\_configs}}.

    \item \textbf{\lr{\texttt{V12\_\_Add\_Course\_Packages.sql}}}:\\
    ساخت جداول فروش بسته‌ای (پکیج جامع چند دوره با تخفیف ویژه) شامل جدول \lr{\texttt{course\_packages}} و جدول پیوندی \lr{\texttt{course\_package\_items}}.

    \item \textbf{\lr{\texttt{V13\_\_Seed\_Sample\_Package.sql}}}:\\
    ایجاد داده‌های پیش‌فرض برای نمونه پکیج جامع تخفیف‌دار دوره‌های آموزشی در کاتالوگ فروش.

    \item \textbf{\lr{\texttt{V14\_\_Add\_Purchases\_Search\_And\_Filtering\_Indexes.sql}}}:\\
    ایجاد ایندکس‌های ترکیبی جهت جستجوی سریع تراکنش‌ها و گزارش‌گیری پرداخت‌ها در پنل مدیریت.

    \item \textbf{\lr{\texttt{V15\_\_Add\_Course\_And\_Package\_ImageUrl.sql}}}:\\
    افزودن ستون آدرس مستقیم تصویر کاور پکیج‌ها و آیکون دوره‌ها در ساختار دیتابیس.

    \item \textbf{\lr{\texttt{V16\_\_Add\_Persistent\_Refresh\_Tokens.sql}}}:\\
    ساخت جدول \lr{\texttt{refresh\_tokens}} جهت ذخیره‌سازی امن و تمدید خودکار نشست احراز هویت کاربران در اپلیکیشن موبایل تا ۳۰ روز.

    \item \textbf{\lr{\texttt{V17\_\_Add\_User\_ProfilePictureUrl.sql}}}:\\
    افزودن فیلد آدرس تصویر پروفایل کاربر (\lr{\texttt{profile\_picture\_url}}) برای بارگذاری عکس آواتار در مسیر uploads.
\end{itemize}

\subsection{دستورات کاربردی اجرای دستی اسکریپت‌های \lr{SQL} روی سرور}
در صورت نیاز به اجرای دستی اسکریپت‌ها یا اتصال مستقیم به محیط خط فرمان پایگاه داده:

\begin{enumerate}
    \item \textbf{ورود به محیط ترمینال دیتابیس (\lr{psql CLI}):}
    \begin{latin}
    \small
    \begin{verbatim}
    docker exec -it leitner-postgres-db psql -U leitner_admin -d leitner_db
    \end{verbatim}
    \end{latin}
    \item \textbf{بررسی جدول نسخه‌های مایگریشن اجرا شده:}
    \begin{latin}
    \small
    \begin{verbatim}
    SELECT scriptname, applied FROM schemaversions ORDER BY applied DESC;
    \end{verbatim}
    \end{latin}
    \item \textbf{اجرای یک فایل \lr{SQL} مشخص بر روی دیتابیس:}
    \begin{latin}
    \small
    \begin{verbatim}
    cat /opt/leitner-platform/deployment/db/migrations/V1__Initial_Schema.sql \
      | docker exec -i leitner-postgres-db psql -U leitner_admin -d leitner_db
    \end{verbatim}
    \end{latin}
\end{enumerate}

% ----------------------------------------------------------------------
% Section 5
% ----------------------------------------------------------------------
\newpage
\section{مشخصات پیشنهادی سرور جدید و پیش‌نیازها}
\begin{itemize}
    \item \textbf{سیستم‌عامل پیشنهادی:} \lr{Ubuntu 24.04 LTS (Noble Numbat)} یا \lr{Ubuntu 22.04 LTS} با معماری \lr{x86\_64}.
    \item \textbf{سخت‌افزار حداقلی:} ۲ هسته پردازنده (CPU)، حداقل ۴ گیگابایت حافظه رم (RAM) و حداقل ۴۰ گیگابایت حافظه \lr{NVMe SSD}.
    \item \textbf{سخت‌افزار پیشنهادی:} ۴ هسته پردازنده، ۸ گیگابایت رم، ۸۰ گیگابایت حافظه \lr{NVMe SSD}.
    \item \textbf{پورت‌های باز فایروال:} پورت ۲۲ (SSH)، پورت ۸۰ (HTTP) و پورت ۴۴۳ (HTTPS).
\end{itemize}

% ----------------------------------------------------------------------
% Section 6
% ----------------------------------------------------------------------
\section{فاز ۱: اقدامات اولیه پیش از روز مهاجرت (۲۴ تا ۴۸ ساعت قبل)}

\subsection{گام ۱.۱: کاهش زمان حیات رکوردهای دی‌ان‌اس (\lr{DNS TTL})}
در پنل مدیریت دی‌ان‌اس دامنه (مانند کلودفلر، ابر آروان یا رجیسترار):
\begin{itemize}
    \item رکورد \lr{A} مربوط به آدرس \verb|api.rightlearn.ir| را بیابید.
    \item رکورد \lr{A} مربوط به آدرس اصلی یا پنل مدیریت (\verb|rightlearn.ir|) را بیابید.
    \item مقدار \lr{TTL} آن‌ها را از حالت خودکار یا ۲۴ ساعت، بر روی \textbf{۳۰۰ ثانیه (۵ دقیقه)} یا \textbf{۶۰ ثانیه} تنظیم نمایید.
    \item \textbf{اهمیت فنی:} این کار موجب می‌شود در زمان تغییر آی‌پی سرور، ترافیک کاربران ظرف چند دقیقه به سرور جدید منتقل گردد و کش شرکت‌های اینترنتی ایجاد تاخیر نکند.
\end{itemize}

\subsection{گام ۱.۲: آماده‌سازی دسترسی کلید امنیتی \lr{SSH}}
کلید عمومی \lr{id\_rsa\_deploy.pub} مربوط به توسعه‌دهنده را روی سرور جدید بارگذاری کنید تا فرآیند ارسال خودکار فایل‌ها و دستورات ترمینال بدون نیاز به وارد کردن مکرر رمز عبور انجام شود:

\begin{latin}
\small
\begin{verbatim}
# From local machine (PowerShell):
type C:\Users\Administrator\.ssh\id_rsa_deploy.pub | ssh root@<NEW_SERVER_IP> \
  "mkdir -p ~/.ssh && cat >> ~/.ssh/authorized_keys && chmod 600 ~/.ssh/authorized_keys"
\end{verbatim}
\end{latin}

% ----------------------------------------------------------------------
% Section 7
% ----------------------------------------------------------------------
\section{فاز ۲: آماده‌سازی و راه‌اندازی زیرساخت در سرور جدید}
با دسترسی \lr{root} از طریق \lr{SSH} به سرور مقصد متصل شوید و مراحل زیر را پیاده‌سازی نمایید:

\subsection{گام ۲.۱: به‌روزرسانی سیستم‌عامل و نصب بسته‌های ضروری}
\begin{latin}
\small
\begin{verbatim}
apt update && apt upgrade -y
apt install -y curl git ufw build-essential unzip tar nginx certbot python3-certbot-nginx
\end{verbatim}
\end{latin}

\subsection{گام ۲.۲: نصب موتور داکر و ابزار \lr{Docker Compose}}
برای نصب مطمئن و پایدار، از مخزن رسمی داکر استفاده نمایید:
\begin{latin}
\small
\begin{verbatim}
# Setup Docker official repository:
install -m 0755 -d /etc/apt/keyrings
curl -fsSL https://download.docker.com/linux/ubuntu/gpg -o /etc/apt/keyrings/docker.asc
chmod a+r /etc/apt/keyrings/docker.asc

echo "deb [arch=$(dpkg --print-architecture) \
  signed-by=/etc/apt/keyrings/docker.asc] \
  https://download.docker.com/linux/ubuntu \
  $(. /etc/os-release && echo "$VERSION_CODENAME") stable" \
  | tee /etc/apt/sources.list.d/docker.list > /dev/null

apt update
apt install -y docker-ce docker-ce-cli containerd.io \
  docker-buildx-plugin docker-compose-plugin

# Verification:
docker --version
docker compose version
\end{verbatim}
\end{latin}

\subsection{گام ۲.۳: پیکربندی دیواره آتش (\lr{UFW Firewall})}
جهت مسدودسازی دسترسی مستقیم خارجی به دیتابیس و پورت‌های داخلی:
\begin{latin}
\small
\begin{verbatim}
ufw default deny incoming
ufw default allow outgoing
ufw allow 22/tcp
ufw allow 80/tcp
ufw allow 443/tcp
ufw --force enable
ufw status
\end{verbatim}
\end{latin}

\subsection{گام ۲.۴: ساخت ساختار پوشه‌های پلتفرم}
\begin{latin}
\small
\begin{verbatim}
mkdir -p /opt/leitner-platform
mkdir -p /var/www/html
\end{verbatim}
\end{latin}

% ----------------------------------------------------------------------
% Section 8
% ----------------------------------------------------------------------
\section{فاز ۳: انتقال اولیه کدها و دارایی‌های حجیم (قبل از پنجره قطعی)}
برای به حداقل رساندن زمان قطعی در لحظه سوییچ نهایی، داده‌های حجیم چندرسانه‌ای را می‌توان زمانی که سرور قدیمی هنوز فعال است انتقال داد:

\subsection{گام ۳.۱: انتقال کدهای اجرایی و فایل‌های پروژه}
با استفاده از اسکریپت استقرار سیستم یا انتقال مستقیم از طریق \lr{SCP}، ساختار پوشه استقرار، باینری‌های بیلد شده و فایل \verb|.env| به سرور جدید منتقل می‌شود:
\begin{latin}
\small
\begin{verbatim}
# Remote project layout on destination server:
/opt/leitner-platform/
  |-- .env
  |-- deployment/
  |-- publish/
  |-- admin-panel/
\end{verbatim}
\end{latin}

\subsection{گام ۳.۲: همگام‌سازی بسته‌های زیپ دوره‌ها و فایل‌های چندرسانه‌ای}
دستورات زیر را روی سرور قدیمی اجرا کنید تا فایل‌های سنگین را به سرور جدید کپی کند:
\begin{latin}
\small
\begin{verbatim}
# On OLD server (45.94.215.188):
rsync -avz /opt/leitner-platform/backend/LeitnerPlatform.API/wwwroot/courses/ \
  root@<NEW_SERVER_IP>:/opt/leitner-platform/backend/LeitnerPlatform.API/wwwroot/courses/

rsync -avz /opt/leitner-platform/backend/LeitnerPlatform.API/wwwroot/uploads/ \
  root@<NEW_SERVER_IP>:/opt/leitner-platform/backend/LeitnerPlatform.API/wwwroot/uploads/
\end{verbatim}
\end{latin}

% ----------------------------------------------------------------------
% Section 9
% ----------------------------------------------------------------------
\section{فاز ۴: آغاز پنجره مهاجرت و انتقال تضمینی دیتابیس (قطع موقت)}
\textbf{نکته حیاتی:} این مرحله باید در زمان کم‌ترافیک (مثلاً ساعات پایانی شب) اجرا گردد. کل این مرحله بین ۵ الی ۱۰ دقیقه طول می‌کشد.

\subsection{گام ۴.۱: قرار دادن سیستم در حالت تعمیرات (\lr{Maintenance Mode})}
برای جلوگیری از ثبت هرگونه سفارش خرید جدید یا مطالعه فلش‌کارت در حین انتقال، سامانه را با اسکریپت پاورشل به حالت تعمیرات سراسری ببرید:
\begin{latin}
\small
\begin{verbatim}
# From local management console:
powershell -ExecutionPolicy Bypass -File .\scripts\manage-admin.ps1 \
  -Action down -Scope Global -ServerIP "45.94.215.188"
\end{verbatim}
\end{latin}
با این کار کاربران اپلیکیشن موبایل بلافاصله صفحه اطلاع‌رسانی تعمیرات را مشاهده خواهند کرد و هیچ تغییری در دیتابیس رخ نمی‌دهد.

\subsection{گام ۴.۲: متوقف کردن کانتینرهای سرویس‌دهنده در سرور قدیمی}
\begin{latin}
\small
\begin{verbatim}
# On OLD server:
cd /opt/leitner-platform
docker compose -f deployment/docker-compose.yml stop backend worker
\end{verbatim}
\end{latin}

\subsection{گام ۴.۳: خروجی کامل (\lr{Dump}) از پایگاه داده لایتنر}
\begin{latin}
\small
\begin{verbatim}
# On OLD server:
docker exec -t leitner-postgres-db pg_dump -U leitner_admin -d leitner_db \
  -F c -b -v -f /tmp/leitner_db_backup.dump

gzip -f /tmp/leitner_db_backup.dump
\end{verbatim}
\end{latin}

\subsection{گام ۴.۴: ارسال فایل پشتیبان به سرور جدید}
\begin{latin}
\small
\begin{verbatim}
# Transfer compressed database dump to NEW server:
scp /tmp/leitner_db_backup.dump.gz root@<NEW_SERVER_IP>:/tmp/
\end{verbatim}
\end{latin}

\subsection{گام ۴.۵: راه‌اندازی کانتینر دیتابیس در سرور جدید}
\begin{latin}
\small
\begin{verbatim}
# On NEW server:
cd /opt/leitner-platform
docker compose -f deployment/docker-compose.yml --env-file .env up -d db redis
\end{verbatim}
\end{latin}
چند ثانیه صبر کنید تا سرویس \verb|PostgreSQL| آماده پذیرش اتصال گردد.

\subsection{گام ۴.۶: بازیابی کامل داده‌ها (\lr{Restore}) در سرور جدید}
\begin{latin}
\small
\begin{verbatim}
# On NEW server:
gunzip -f /tmp/leitner_db_backup.dump.gz

# Drop existing blank schema and create target database:
docker exec -i leitner-postgres-db psql -U leitner_admin -d postgres \
  -c "DROP DATABASE IF EXISTS leitner_db;"
docker exec -i leitner-postgres-db psql -U leitner_admin -d postgres \
  -c "CREATE DATABASE leitner_db OWNER leitner_admin;"

# Copy and restore database dump:
docker cp /tmp/leitner_db_backup.dump leitner-postgres-db:/tmp/leitner_db_backup.dump
docker exec -i leitner-postgres-db pg_restore -U leitner_admin -d leitner_db \
  -v /tmp/leitner_db_backup.dump

# Clean temporary files:
docker exec leitner-postgres-db rm -f /tmp/leitner_db_backup.dump
rm -f /tmp/leitner_db_backup.dump
\end{verbatim}
\end{latin}

% ----------------------------------------------------------------------
% Section 10
% ----------------------------------------------------------------------
\section{فاز ۵: راه‌اندازی کامل کانتینرها در سرور جدید}
اکنون تمام سرویس‌های پلتفرم را در سرور مقصد استارت کنید:
\begin{latin}
\small
\begin{verbatim}
# On NEW server:
cd /opt/leitner-platform
docker compose -f deployment/docker-compose.yml --env-file .env up -d --build

# Verify all containers are running:
docker compose -f deployment/docker-compose.yml ps
\end{verbatim}
\end{latin}
باید وضعیت هر ۵ کانتینر در حالت \lr{Up} نمایش داده شود:
\begin{itemize}
    \item \lr{leitner-postgres-db}
    \item \lr{leitner-redis-cache}
    \item \lr{leitner-backend-api}
    \item \lr{leitner-admin-panel}
    \item \lr{leitner-background-worker}
\end{itemize}

% ----------------------------------------------------------------------
% Section 11
% ----------------------------------------------------------------------
\section{فاز ۶: پیکربندی پروکسی معکوس \lr{Nginx} و گواهی امنیتی \lr{SSL}}

\subsection{گام ۶.۱: انتقال صفحه پشتیبان تعمیرات}
\begin{latin}
\small
\begin{verbatim}
cp /opt/leitner-platform/deployment/maintenance.html /var/www/html/maintenance.html
\end{verbatim}
\end{latin}

\subsection{گام ۶.۲: ساخت کانفیگ پروکسی در \lr{Nginx}}
فایل \verb|/etc/nginx/sites-available/leitner.conf| را با محتوای زیر روی سرور جدید ایجاد کنید:

\begin{latin}
\footnotesize
\begin{verbatim}
# 1. API Reverse Proxy for Mobile App
server {
    listen 80;
    server_name api.rightlearn.ir;
    client_max_body_size 200M;

    location / {
        proxy_pass http://127.0.0.1:8080;
        proxy_http_version 1.1;
        proxy_set_header Upgrade $http_upgrade;
        proxy_set_header Connection keep-alive;
        proxy_set_header Host $host;
        proxy_cache_bypass $http_upgrade;
        proxy_set_header X-Real-IP $remote_addr;
        proxy_set_header X-Forwarded-For $proxy_add_x_forwarded_for;
        proxy_set_header X-Forwarded-Proto $scheme;
        proxy_read_timeout 300s;
        proxy_connect_timeout 300s;
        proxy_send_timeout 300s;
    }
}

# 2. Admin Panel & Web Gateway
server {
    listen 80;
    server_name rightlearn.ir admin.rightlearn.ir;
    client_max_body_size 200M;

    error_page 502 503 504 /maintenance.html;
    location = /maintenance.html {
        root /var/www/html;
        internal;
    }

    location / {
        proxy_pass http://127.0.0.1:8081;
        proxy_intercept_errors on;
        proxy_http_version 1.1;
        proxy_set_header Upgrade $http_upgrade;
        proxy_set_header Connection keep-alive;
        proxy_set_header Host $host;
        proxy_cache_bypass $http_upgrade;
        proxy_set_header X-Real-IP $remote_addr;
        proxy_set_header X-Forwarded-For $proxy_add_x_forwarded_for;
        proxy_set_header X-Forwarded-Proto $scheme;
        proxy_read_timeout 300s;
        proxy_connect_timeout 300s;
        proxy_send_timeout 300s;
    }
}
\end{verbatim}
\end{latin}

فعال‌سازی و بازخوانی \lr{Nginx}:
\begin{latin}
\small
\begin{verbatim}
ln -sf /etc/nginx/sites-available/leitner.conf /etc/nginx/sites-enabled/
nginx -t
systemctl reload nginx
\end{verbatim}
\end{latin}

% ----------------------------------------------------------------------
% Section 12
% ----------------------------------------------------------------------
\section{فاز ۷: تغییر رکوردهای دی‌ان‌اس و صدور گواهی \lr{SSL}}

\subsection{گام ۷.۱: تغییر نشانی آی‌پی در پنل دی‌ان‌اس (\lr{DNS Cutover})}
در پنل دامنه خود (کلودفلر / ابر آروان):
\begin{itemize}
    \item آی‌پی رکورد \lr{A} دامنه \verb|api.rightlearn.ir| را به \textbf{آی‌پی سرور جدید} تغییر دهید.
    \item آی‌پی رکورد \lr{A} دامنه \verb|rightlearn.ir| را به \textbf{آی‌پی سرور جدید} تغییر دهید.
\end{itemize}

\subsection{گام ۷.۲: صدور خودکار گواهی \lr{SSL} با \lr{Certbot}}
پس از انتشار دی‌ان‌اس (حدود ۱ الی ۵ دقیقه):
\begin{latin}
\small
\begin{verbatim}
# On NEW server:
certbot --nginx -d api.rightlearn.ir -d rightlearn.ir
\end{verbatim}
\end{latin}
ابزار \lr{Certbot} به صورت خودکار ریدایرکت \lr{HTTPS} روی پورت ۴۴۳ را فعال ساخته و زمان‌بندی تمدید دوره‌ای آن را ثبت می‌کند.

\subsection{گام ۷.۳: خروج از حالت تعمیرات و فعال‌سازی سرویس‌دهی}
با دستور زیر روی سرور جدید، حالت تعمیرات را غیرفعال کنید:
\begin{latin}
\small
\begin{verbatim}
# On NEW server:
echo "UPDATE system_configs SET value = 'false' WHERE key = 'maintenance_mode';" \
  | docker exec -i leitner-postgres-db psql -U leitner_admin -d leitner_db
\end{verbatim}
\end{latin}
اکنون اپلیکیشن‌های موبایل کاربران و پنل مدیریت مجدداً به صورت عادی و با اتصال به سرور جدید به کار خود ادامه می‌دهند.

% ----------------------------------------------------------------------
% Section 13
% ----------------------------------------------------------------------
\section{فاز ۸: به‌روزرسانی اسکریپت‌ها و پایپ‌لاین‌های استقرار}
برای آنکه در بیلدهای بعدی فایل‌ها به سرور جدید ارسال شوند، موارد زیر در مخزن پروژه به‌روزرسانی می‌شوند:

\begin{enumerate}
    \item \textbf{اسکریپت استقرار:} مقدار متغیر \verb|$ServerIP| در فایل \verb|scripts/deploy-to-server.ps1| به آی‌پی سرور جدید تغییر می‌یابد.
    \item \textbf{اسکریپت مدیریت لایف‌سایکل:} مقدار متغیر \verb|$ServerIP| در فایل \verb|scripts/manage-admin.ps1| به آی‌پی سرور جدید تغییر می‌یابد.
    \item \textbf{پایپ‌لاین گیت‌هاب اکشنز (\lr{GitHub Actions}):}
    \begin{itemize}
        \item در مسیر مخزن گیت‌هاب به بخش \lr{Settings > Secrets and variables > Actions} مراجعه فرمایید.
        \item متغیر مخفی \verb|SERVER_IP| را به آی‌پی سرور جدید تغییر دهید.
        \item در صورتی که کلید \lr{SSH} جدیدی تولید کرده‌اید، مقدار \verb|SERVER_SSH_KEY| را با کلید خصوصی جدید به‌روزرسانی کنید.
    \end{itemize}
\end{enumerate}

% ----------------------------------------------------------------------
% Section 14
% ----------------------------------------------------------------------
\section{فاز ۹: چک‌لیست اعتبارسنجی و تست‌های پذیرش فنی}
پس از اتمام مهاجرت، آزمون‌های جدول زیر توسط تیم تست و کارفرما انجام می‌پذیرد:

\begin{table}[ht]
\centering
\caption{چک‌لیست تست‌های صحت عملکرد سامانه پس از مهاجرت}
\vspace{2mm}
\begin{tabularx}{\textwidth}{lXX}
\toprule
\textbf{عنوان آزمون} & \textbf{دستور یا اقدام بررسی} & \textbf{نتیجه مورد انتظار} \\
\midrule
سلامت کانتینرها & \lr{docker compose ps} & تمام ۵ کانتینر در وضعیت Up هستند \\
پاسخ‌دهی وب‌سرویس & درخواست به \lr{api/v1/auth/captcha} & کد وضعیت 200 به همراه داده تصویر امنیتی \\
اتصال دیتابیس & لاگ کانتینر \lr{leitner-backend-api} & اعمال موفق مایگریشن‌ها و شنود پورت 8080 \\
پنل مدیریت وب & باز کردن \lr{https://rightlearn.ir} & نمایش صفحه ورود با ظاهر مناسب \\
دانلود دوره‌ها & دانلود یک فایل از مسیر \lr{/courses/} & دانلود کامل فایل zip بدون خطا \\
عملکرد اپلیکیشن & ورود به اپ و مطالعه کارت‌ها & همگام‌سازی سریع بدون خطای اتصال سرور \\
سرویس پیامک & ارسال کد یکبارمصرف (OTP) & دریافت آنی پیامک تایید هویت \\
\bottomrule
\end{tabularx}
\end{table}

% ----------------------------------------------------------------------
% Section 15
% ----------------------------------------------------------------------
\section{طرح بازگشت به شرایط قبل (\lr{Rollback Plan}) و خروج سرور قدیم}

\subsection{سناریوی بازگشت به شرایط قبل در صورت بروز مشکل غیرمنتظره}
اگر در هر یک از مراحل راه‌اندازی سرور جدید با مشکل سخت‌افزاری یا شبکه مواجه شدید، به راحتی ظرف ۲ دقیقه می‌توان به سرور قبلی بازگشت:
\begin{enumerate}
    \item رکوردهای \lr{A} در پنل \lr{DNS} را به آی‌پی سرور قبلی (\verb|45.94.215.188|) برگردانید.
    \item در سرور قبلی، کانتینرهای بک‌اند را روشن نموده و حالت تعمیرات را خاموش کنید:
    \begin{latin}
    \small
    \begin{verbatim}
    docker compose -f deployment/docker-compose.yml start backend worker
    echo "UPDATE system_configs SET value = 'false' WHERE key = 'maintenance_mode';" \
      | docker exec -i leitner-postgres-db psql -U leitner_admin -d leitner_db
    \end{verbatim}
    \end{latin}
\end{enumerate}

\subsection{از مدار خارج کردن سرور قدیمی (\lr{Decommissioning})}
پس از گذشت ۲۴ الی ۴۸ ساعت از عملکرد بی‌نقص سرور جدید:
\begin{enumerate}
    \item در سرور قدیمی، تمام کانتینرها متوقف و سرویس \lr{Nginx} خاموش شود:
    \begin{latin}
    \small
    \begin{verbatim}
    cd /opt/leitner-platform && docker compose -f deployment/docker-compose.yml down
    systemctl stop nginx && systemctl disable nginx
    \end{verbatim}
    \end{latin}
    \item سرور قدیمی از پنل ارائه‌دهنده سرویس ابری حذف یا خاموش گردد.
\end{enumerate}

\end{document}
